% sections/market-overview.tex
\section{Market Overview}

\subsection{Category size and growth}

Global spend on expense management and corporate card software is estimated at \$9.4B for
FY2026, growing 19\% year over year as finance teams replace manual expense reports and
plastic-only corporate cards with unified platforms \cite{cedarpoint2026sizing}. Growth is
concentrated in two motions: mid-market companies (200--2,000 employees) buying their first
dedicated platform, and enterprise accounts consolidating multiple point tools -- a separate
card issuer, a separate expense tool, a separate travel booking system -- onto one vendor.

\begin{center}
\footnotesize
\begin{tabularx}{\linewidth}{@{}Y r r@{}}
\toprule
\rowcolor{ink}
\textcolor{white}{\textbf{Segment}} & \textcolor{white}{\textbf{FY2026 share}} & \textcolor{white}{\textbf{CAGR (FY2026--FY2029)}} \\
\midrule
Mid-market (200--2,000 employees) & 44\% & 22\% \\
\rowcolor{mist}
Enterprise (2,000+ employees) & 38\% & 17\% \\
SMB (under 200 employees) & 18\% & 13\% \\
\bottomrule
\end{tabularx}
\end{center}

\subsection{Buying committee and decision drivers}

Deals are rarely won on a single feature. Across 214 competitive deals logged in the last four
quarters, the buying committee consistently weighs the same four factors, in this order:

\begin{enumerate}[leftmargin=1.4em, itemsep=2pt]
  \item \textbf{Control and visibility} -- real-time card limits, policy enforcement at the
    point of spend, and consolidated reporting across entities and currencies.
  \item \textbf{Time to value} -- weeks versus months to go live, and how much finance-team
    time an implementation consumes before the first close cycle runs through the new system.
  \item \textbf{Integration depth} -- native, bidirectional sync with the general ledger and
    ERP, not a nightly CSV export.
  \item \textbf{Total cost of ownership} -- subscription price, card interchange economics, and
    the hidden cost of professional-services hours required to keep the platform configured.
\end{enumerate}

\begin{tcolorbox}[infobox, title={\faChartBar\ Category shift to note}]
The market is moving from best-of-breed toward consolidated platforms. 61\% of enterprise RFPs
issued in the last two quarters explicitly asked for expense, card issuing, and at least
lightweight procurement in a single contract -- up from 38\% a year earlier
\cite{almanac2026procurement}. Vendors that stay expense-only are increasingly evaluated as a
point solution rather than a platform, which shows up directly in win rates above 2,000
employees.
\end{tcolorbox}

\subsection{Where Nexa plays}

Nexa's installed base is concentrated in mid-market technology, professional services, and
healthcare companies in North America and the United Kingdom, with average contract value of
\$38,000 and a median sales cycle of 71 days. We have almost no presence in continental Europe,
where VAT reclaim and multi-currency consolidation requirements favor Meridian and Vertex, and
limited traction above 3,000 employees, where multi-entity consolidation is table stakes.
